\section{Cocycle-twisted regular representations}
\label{sec:projective_twisted_regular}

\begin{definition}[Cocycle-twisted regular matrices]
    \label{def:mpug_twisted_regular_matrices}
    \lean{ScalarCocycle.phase,
        ScalarCocycle.twistedLeftRegular,
        ScalarCocycle.twistedRightRegular}
    \leanok
    \uses{def:scalar_cocycle, thm:mpug_circle_complex_units_h2}
    Let $G$ be a group and let $\omega$ be a $\C^\times$-valued scalar
    cochain.  Write
    \begin{align}
        \phi(g,h) & =\frac{\omega(g,h)}{|\omega(g,h)|}\in U(1).
    \end{align}
    On the standard basis of $\C^G$, define
    \begin{align}
        L_g^\phi
         & =\sum_{h\in G}\phi(g,h)^{-1}|gh\rangle\langle h|,
         &
        R_g^\phi
         & =\sum_{h\in G}\phi(h,g^{-1})^{-1}|hg^{-1}\rangle\langle h|.
        \label{eq:mpug_twisted_regular_matrices}
    \end{align}
    Whereas \cite[lines~412--418]{FrancoRubio2025MPUGauging} starts from a
    $U(1)$-valued cocycle $\Omega$, taking $\Omega=\phi$ reproduces its
    inverse conventions.
\end{definition}

\begin{lemma}[Cocycle equation for the circle phase]
    \label{lem:mpug_circle_phase_cocycle}
    \lean{ScalarCocycle.phase_isCocycle}
    \leanok
    \uses{def:is_cocycle, thm:mpug_circle_complex_units_h2}
    If $\omega$ is a cocycle, then its circle phase satisfies
    \begin{align}
        \phi(g,h)\phi(gh,k)
         & =\phi(g,hk)\phi(h,k)
    \end{align}
    for all $g,h,k\in G$.
\end{lemma}

\begin{proof}\leanok
    \uses{thm:mpug_circle_complex_units_h2}
    The pointwise phase map $z\mapsto z/|z|$ is multiplicative, so it preserves
    the cocycle equation.
\end{proof}

\begin{lemma}[Action of the twisted regular matrices]
    \label{thm:mpug_twisted_regular_action}
    \lean{ScalarCocycle.twistedLeftRegular_mulVec_apply,
        ScalarCocycle.twistedRightRegular_mulVec_apply}
    \leanok
    \uses{def:mpug_twisted_regular_matrices}
    If $G$ is finite, then for $v\in\C^G$ and $g,h\in G$,
    \begin{align}
        (L_g^\phi v)(gh)
         & =\phi(g,h)^{-1}v(h),
         &
        (R_g^\phi v)(hg^{-1})
         & =\phi(h,g^{-1})^{-1}v(h).
    \end{align}
    % Source anchor: arXiv:2502.20257, lines 412--418.
\end{lemma}

\begin{proof}\leanok
    \uses{def:mpug_twisted_regular_matrices}
    In each matrix-vector product, exactly one summand is nonzero: the column
    indexed by $h$.
\end{proof}

\begin{lemma}[Multiplication laws for twisted regular matrices]
    \label{thm:mpug_twisted_regular_multiplication}
    \lean{ScalarCocycle.twistedLeftRegular_mul,
        ScalarCocycle.twistedRightRegular_mul}
    \leanok
    \uses{def:mpug_twisted_regular_matrices,
        lem:mpug_circle_phase_cocycle}
    If $G$ is finite and $\omega$ is a cocycle, then
    \begin{align}
        L_g^\phi L_h^\phi
         & =\phi(g,h)^{-1}L_{gh}^\phi,           \\
        R_g^\phi R_h^\phi
         & =\phi(h^{-1},g^{-1})^{-1}R_{gh}^\phi.
        \label{eq:mpug_twisted_regular_multiplication}
    \end{align}
    % Source anchor: arXiv:2502.20257, lines 406--418.
\end{lemma}

\begin{proof}\leanok
    \uses{lem:mpug_circle_phase_cocycle}
    Expanding the matrix products entrywise reduces the two scalar identities
    to the cocycle equation at $(g,h,k)$ and $(k,h^{-1},g^{-1})$, respectively.
\end{proof}

\begin{lemma}[Commutation of twisted left and right actions]
    \label{thm:mpug_twisted_regular_commutation}
    \lean{ScalarCocycle.twistedLeftRegular_commute}
    \leanok
    \uses{def:mpug_twisted_regular_matrices}
    If $G$ is finite and $\omega$ is a cocycle, then
    \begin{align}
        L_g^\phi R_h^\phi
         & =R_h^\phi L_g^\phi.
        \label{eq:mpug_twisted_regular_commutation}
    \end{align}
    % Source anchor: arXiv:2502.20257, lines 406--418 and 457.
\end{lemma}

\begin{proof}\leanok
    \uses{lem:mpug_circle_phase_cocycle}
    Expanding both matrix products entrywise and applying the cocycle equation
    at $(g,k,h^{-1})$ yields the equality.
\end{proof}

\begin{lemma}[Unitarity of twisted regular matrices]
    \label{thm:mpug_twisted_regular_unitary}
    \lean{ScalarCocycle.twistedLeftRegular_mem_unitaryGroup,
        ScalarCocycle.twistedRightRegular_mem_unitaryGroup}
    \leanok
    \uses{def:mpug_twisted_regular_matrices}
    If $G$ is finite, the matrices $L_g^\phi$ and $R_g^\phi$ are unitary for
    every $g\in G$.
\end{lemma}

\begin{proof}\leanok
    \uses{def:mpug_twisted_regular_matrices}
    Each matrix is a permutation matrix whose nonzero entries have modulus
    one.  Distinct columns therefore have disjoint support, while every column
    has squared norm one.
\end{proof}

\begin{lemma}[A coefficient identity for the circle phase]
    \label{thm:mpug_onsite_gauss_phase}
    \lean{ScalarCocycle.onsiteGauss_phase_div}
    \leanok
    \uses{def:mpug_twisted_regular_matrices,
        lem:mpug_circle_phase_cocycle}
    Suppose that $\omega$ is a cocycle and impose the additional normalization
    \begin{align}
        \phi(t^{-1},t) & =1
    \end{align}
    for every $t\in G$.  Then
    \begin{align}
        \frac{\phi(ag^{-1},gb)}{\phi(a,b)}
         & =\frac{1}{\phi(a,g^{-1})\phi(g,b)}.
        \label{eq:mpug_onsite_gauss_phase}
    \end{align}
    The scalar convention $S_{t^{-1}}=S_t^{-1}$ appears in
    \cite[line~238]{FrancoRubio2025MPUGauging}, and the coefficient rewrite is
    at lines 3763--3766. The displayed normalization remains an explicit
    assumption here; it is not a consequence of cocyclicity.
\end{lemma}

\begin{proof}\leanok
    \uses{lem:mpug_circle_phase_cocycle}
    Applying the cocycle equation first at $(a,g^{-1},gb)$ and then at
    $(g^{-1},g,b)$, together with $\phi(e,b)=1$ from the normalization at
    $t=e$, yields the identity.
\end{proof}

\begin{theorem}[A unitary projective representation for every class]
    \label{thm:twisted_regular_projective}
    \lean{ScalarCocycle.inverse,
        ScalarCocycle.inverse_isCocycle,
        ScalarCocycle.inverse_circlePhaseInclusion,
        ScalarCocycle.twistedLeftRegularGL,
        ScalarCocycle.twistedLeftRegularGL_mul,
        ScalarCocycle.twistedLeftRegularProjective,
        ScalarCocycle.twistedLeftRegularProjective_mem_unitaryGroup,
        ScalarCocycle.regularProjectiveForClass,
        ScalarCocycle.regularProjectiveForClass_mem_unitaryGroup}
    \leanok
    \uses{def:projective_rep, thm:mpug_twisted_regular_multiplication,
        thm:mpug_twisted_regular_unitary, thm:mpug_circle_complex_units_h2}
    Let $G$ be a finite group and $\omega$ a $\C^\times$-valued $2$-cocycle.
    There is a unitary projective representation of $G$ on $\C^G$ whose
    factor system is the pointwise phase $\omega/|\omega|$, a cocycle
    cohomologous to $\omega$. This is the input of the construction
    \cite[Section~III.A]{Cirac2021Matrix} starting from any solution of the
    cocycle condition.
\end{theorem}

\begin{proof}\leanok
    \uses{thm:mpug_twisted_regular_multiplication,
        thm:mpug_twisted_regular_unitary, thm:mpug_circle_complex_units_h2}
    The pointwise inverse $\omega^{-1}$ is again a cocycle, and its pointwise
    phase is the inverse of the phase $\bar\omega=\omega/|\omega|$ of
    $\omega$. By Lemma~\ref{thm:mpug_twisted_regular_multiplication}, the
    twisted left regular matrices $L_g$ of $\omega^{-1}$ satisfy
    $L_gL_h=\bar\omega(g,h)L_{gh}$, and by
    Lemma~\ref{thm:mpug_twisted_regular_unitary} they are unitary, hence
    invertible. Relabelling $G$ by $\{1,\dots,|G|\}$ gives the representation.
    By Lemma~\ref{thm:mpug_circle_complex_units_h2}, $\bar\omega$ is
    cohomologous to $\omega$ because $G$ is finite.
\end{proof}
